\cotefile{156-6}
\scanpage{42}\mspage{\deleted{\unsure{39|34}} 41}

\deleted{le sous-atelier $\unsure{\mathcal F|\mathcal F'}$ à partir de $\mathcal F$, comme} \deleted{\unsure{l'ens des} $F$}
\[
\deleted{\mathcal F'=\{F\in\mathcal F\mid\widetilde F\subset M'\}\mathrel{\overset{\mathrm{dfn}}{=}}\mathrm{Env}(M')}
\]
\deleted{Donc} \deleted{On trouve ainsi}

\uline{Proposition} (Ne dépend pas en fait \unsure{d'une} donnée $\ll$ sur $\mathcal F$, \inserted{seulemt de $\leq$}) On a deux applications \unsure{biunivoques} inverses l'une de l'autre
\diagram{\begin{tikzcd}[ampersand replacement=\&,column sep=large,row sep=small]
{\substack{\deleted{\unsure{\mathrm{Satsp}|\mathrm{Satpr}}}\\\mathrm{Sousatspac}(\mathcal F)}} \arrow[r,shift left=1ex,"{\mathcal F'\longmapsto\mathcal F'\cap M}"] \& \mathrm{Ferm}(M) \arrow[l,shift left=1ex] \\
{\{F\in\mathcal F\mid\widetilde F\subset M\}\mathrel{\overset{\mathrm{dfn}}{=}}\mathrm{Env}(M')} \& M' \arrow[l,maps to]
\end{tikzcd}}
% slanted
\margin{\deleted{\unsure{Dans le cas} $\mathcal F$} Si on identifie $\mathcal F$ à une partie de $\unsure{\mathcal P|\mathrm P}(M)$, alors $\mathrm{Env}.(M')=\{F\in\mathcal F\mid F\subset M'$}
entre l'ensemble des \deleted{\unsure{sous-ateliers}} \inserted{parties \unsure{propres|fermées}} spacieuses \inserted{$\mathcal F'$} de $\mathcal F$, et \unsure{l'ens} des parties \uline{fermées} \inserted{$M'$} de $M$.

Il faut voir, \deleted{que si} \inserted{quand $\mathcal F$ est un atelier, si} $\mathcal F'$ est un sous-atelier de $\mathcal F$, \deleted{\unsure{alors}} pour $\leq$, $\ll$, \deleted{\unsure{c'est un}} \inserted{au sens des} ateliers $\mathrm{At}\,1$ -- $\mathrm{At}\,7$. Rappelons pour mémoire $\mathrm{At}\,1$ -- $\mathrm{At}\,3$

\uline{At 1} \deleted{\unsure{Les parties} \illegible{1}} \unsure{Les} sup majorés \inserted{(pour $\leq$)} \unsure{existent} dans $\mathcal F$ ; et $\mathcal F\ne\varnothing$ (ou encore, a un plus petit élément $\varnothing_{\mathcal F}$).

\uline{At} 2 Toute \deleted{\unsure{fig}} $F\in\mathcal F$ est borne sup. pour $\leq$ des ``figures élémentaires'' (\unsure{les|des} \unsure{pts|élts|elts} \unsure{strictmt.} irr.)\unsure{.} qu'elle majore (pour $\leq$)

\uline{At} 3 Si $F,G\in\mathcal F$, pour que $F\vee\unsure{C|G}$ existe dans $\mathcal F$, \inserted{(\unsure{i.e.|i.e.,} $\{F,G\}$ majoré)} il suffit que \deleted{\unsure{si}} pour $X\mathrel{\triangleleft}F$, $Y\mathrel{\triangleleft}C$, $X\vee Y$ existe.
